\section{The non-induced Hecke cutoff obstruction and its literal smooth domain}
\label{sec:hecke-cutoff-obstruction}

The universality input missing from the scope of MA-MO-002 is
available for an imaginary quadratic field. A second change is
favorable: a nontrivial finite-order Hecke function is entire, so
the original ideal sum has no polar discrepancy on the exponential
Schwartz space. We prove both statements with the actual domains
specified. The compact-character identification is the one already
proved in MA-CC-001 and MT-M24; no radial section for an arbitrary
ideal class is assumed.

Let \(K\) be imaginary quadratic and let \(\eta\) be a fixed
nontrivial primitive finite-order Hecke character. The character
need not be induced from \(\mathbb Q\). Extend its coefficients
by zero at conductor ideals as in (MA-RS.40). For a finite set
\(S\) of prime ideals put
\[
 \begin{aligned}
 A(t)&=E_{\eta,S}(\sigma+it)L_\eta(\sigma+it),
       &j(t)&=A(t)^{-1},\\
 E_{\eta,S}(s)&=\prod_{\mathfrak p\in S}
             (1-\eta(\mathfrak p)(N\mathfrak p)^{-s}),
       &a&=\sigma-\tfrac12,\qquad \tfrac12<\sigma<1.
 \end{aligned}
 \tag{MA-MO.100}
\]
Use \(\mathcal H_a,V_a,T_\phi\) from (MA-MO.20)--(MA-MO.22),
where \(\phi\) is any nonconstant bounded smooth function with
compactly supported derivative. The closed operators
\(Z_{a,\eta,S}\) and \(J_{a,\eta,S}\) are exactly (MA-RS.41).
Their compact-frequency core for the present composition is
\[
 \mathscr C_{\eta,\sigma}
 =C_c^\infty\bigl(\mathbb R\setminus
                  \{t:L_\eta(\sigma+it)=0\}\bigr).
 \tag{MA-MO.101}
\]
Write \(\mathcal E\) for the exponential Schwartz space and
\(\mathcal I_\eta\) for the uncentered closed ideal of MT-M18.

\paragraph{MA-MO-009 (unbounded Hecke cutoff on actual smooth ideal sums).}
The composition
\[
 Q_{\eta,\phi,S}=Z_{a,\eta,S}M_\phi J_{a,\eta,S}
 \tag{MA-MO.102}
\]
is defined on \(V_a^{-1}\mathscr C_{\eta,\sigma}\).
Neither it nor \(M_\phi-Q_{\eta,\phi,S}\) has a bounded
extension to \(\mathcal H_a\). This assertion holds for every
line \(1/2<\sigma<1\), including a line with zeros or a line
proved zero free.

More strongly, for every \(F\in\mathcal E\) the literal sum
\[
 (Z_{\eta,S}^{\rm geom}F)(x)
   =\sum_{\substack{\mathfrak a\text{ integral}\\
                         (\mathfrak a,S)=1}}
       \eta(\mathfrak a)(N\mathfrak a)^{-1/2}
                              F(x+\log N\mathfrak a)
 \tag{MA-MO.103}
\]
converges with all derivatives locally uniformly and belongs to
the actual centered smooth theta ideal \(T_{1/2}\mathcal I_\eta\).
For \(y=Z_{\eta,S}^{\rm geom}F\), the original wide-weight
inverse \(J_{\eta,S}^{\rm geom}\) and the narrow maximal inverse
both return \(F\), and
\[
 \begin{aligned}
 Z_{a,\eta,S}F&=Z_{\eta,S}^{\rm geom}F=y,\\
 J_{a,\eta,S}y&=J_{\eta,S}^{\rm geom}y=F,\\
 Q_{\eta,\phi,S}y
   &=Z_{\eta,S}^{\rm geom}(\phi F)
                    \in T_{1/2}\mathcal I_\eta,\\
 M_\phi y-Q_{\eta,\phi,S}y&\in\mathcal E.
 \end{aligned}
 \tag{MA-MO.104}
\]
There is a sequence of such genuine smooth vectors \(y_M\)
with bounded \(\mathcal H_a\) norm for which the norms of
both last two outputs in (MA-MO.104) tend to infinity.
Thus the same unboundedness holds for the literal ideal-sum
conjugation on this specified smooth domain, without a polar
subtraction or an output outside the Hilbert space.

The conclusion applies in particular to the non-induced cubic
character of \(\mathbb Q(\zeta_3)\) in MA-CH-003, and to its
conjugate. The smooth-domain argument also applies to every
nontrivial primitive Dirichlet character, using MA-MO-002 for
its universality step. No use of \(A\) or \(M\) enters these
assertions.

\begin{proof}
\emph{The precise additional published input.}
Mishou (2003), the unnumbered main Theorem on printed p.~46,
proves universality on compact sets with connected complement
in \(\sigma_K<\operatorname{Re}s<1\), where
\(\sigma_K=1-1/[K:\mathbb Q]\) for \(K\ne\mathbb Q\).
It applies to a Gr\"ossencharacter as defined in \S2,
printed pp.~47--49. Our finite-order character has trivial
archimedean type and satisfies that definition. Since
\([K:\mathbb Q]=2\), its permitted strip is exactly
\(1/2<\operatorname{Re}s<1\). The target must be continuous
and nonzero on the compact set and holomorphic inside. The
theorem gives positive lower density of approximating shifts;
only existence is used here. There is no zero-location or
extra zero-density hypothesis.

\emph{Actual composition and the norm obstruction.}
The zero set removed in (MA-MO.101) is discrete, since the
analytic function is not identically zero. Its complement
supports a dense compact smooth core in \(L^2\). On this
core \(jv\) is compact smooth. Multiplication by \(\phi\)
preserves Schwartz space, so \(T_\phi(jv)\) is Schwartz.
The fixed-strip polynomial bounds for \(L_\eta\) and all
its derivatives, already supplied by NT-Z10/MT-M18 and used
in MA-RS-005, imply \(AT_\phi(jv)\in L^2\). Thus every
sequential domain needed in (MA-MO.102) is verified.

Here are the full changes in the norm argument of MA-MO-002.
Set \(h=\sigma-3/4\) and choose \(|h|<r<1/4\).
The compact disk \(C=\{3/4+z:|z|\le r\}\) lies in the
permitted strip and has connected complement. With
\(q(v)=i\widehat{\phi'}_+(v)/(2\pi v)\), choose a fixed
\(d>0\) so small that \(h^2+d^2<r^2\) and \(q(d)\ne0\).
Such a choice follows from the identity theorem for the
nonzero entire transform of \(\phi'\).

For each positive integer \(M\), the target
\(f_M(3/4+z)=e^{-iM(z-h)}\) is entire and nonvanishing,
with minimum modulus \(e^{-Mr}\) on \(C\).
Apply Mishou's theorem with tolerance
\(\epsilon e^{-Mr}\), where \(0<\epsilon<1/100\).
It gives \(T_M\) with
\[
 L_\eta(3/4+z+iT_M)
   =e^{-iM(z-h)}(1+\delta_M(z)),\qquad
      |\delta_M(z)|<\epsilon\quad(|z|\le r).
 \tag{MA-MO.105}
\]
Choose a nonnegative real \(b\in C_c^\infty((-1,1))\)
of \(L^2\) norm one, and use the two unit vectors
\[
 v_M(u)=M^{1/2}b(M(u-T_M)),\qquad
 w_M(t)=M^{1/2}b(M(t-T_M-d)).
 \tag{MA-MO.106}
\]
For large \(M\) their disjoint windows lie inside the
disk and contain no zero, by (MA-MO.105). Off the diagonal
the kernel of the already defined operator with \(S\) empty
is exactly
\(L_\eta(\sigma+it)q(t-u)/L_\eta(\sigma+iu)\).
After \(u=T_M+\xi/M\), \(t=T_M+d+\nu/M\), its pairing
with these vectors is
\[
 \frac{e^{Md}}{M}\iint b(\xi)b(\nu)e^{\nu-\xi}
 q\!\left(d+\frac{\nu-\xi}{M}\right)
 \frac{1+\delta_M(h+i(d+\nu/M))}
      {1+\delta_M(h+i\xi/M)}\,d\xi\,d\nu.
 \tag{MA-MO.107}
\]
The relative error ratio differs from one by at most
\(2\epsilon/(1-\epsilon)\). Uniform continuity of \(q\)
on these windows makes its ratio to \(q(d)\) differ from
one by less than \(1/100\) for large \(M\). The combined
relative error is less than \(1/2\). The remaining weight
is nonnegative and has the strictly positive integral
\(I_b=(\int b(\nu)e^\nu d\nu)(\int b(\xi)e^{-\xi}d\xi)\).
Consequently the pairing has modulus at least
\(\tfrac12|q(d)|I_b e^{Md}/M\). This is a norm bound
on actual vectors, not an inference from a kernel value.

Each finite prime-ideal factor and its inverse is bounded
on this line: use \(1-(N\mathfrak p)^{-\sigma}>0\)
and \(1+(N\mathfrak p)^{-\sigma}\). The exact identity
\[
 M_A T_\phi M_j
  =M_{E_{\eta,S}}
       (M_{L_\eta}T_\phi M_{L_\eta^{-1}})
                               M_{E_{\eta,S}}^{-1}
 \tag{MA-MO.108}
\]
therefore preserves unboundedness, using normalized inputs
\(E_{\eta,S}v_M/\|E_{\eta,S}v_M\|_2\). Subtraction of
the bounded \(T_\phi\), and conjugation by \(V_a\), prove
the first assertion.

\emph{Literal ideal sums and the absence of a polar term.}
For \(F\in\mathcal E\), every derivative of the series
(MA-MO.103) converges absolutely and uniformly for \(x\)
in a compact interval. Indeed exponential decay at the
right bounds its terms by a constant times
\((N\mathfrak a)^{-c-1/2}\), for any fixed \(c>1/2\),
and that ideal series converges. On \(\mathcal H_b\),
\(b>1/2\), it is exactly the absolutely norm-convergent
operator of (MA-CC.1), with Mellin multiplier
\(E_{\eta,S}(z+1/2)L_\eta(z+1/2)\) on
\(\operatorname{Re}z=b\).

The function
\[
 H_F(z)=E_{\eta,S}(z+1/2)L_\eta(z+1/2)\mathcal MF(z)
 \tag{MA-MO.109}
\]
is entire. This is precisely where nontriviality of
\(\eta\) is used: the finite Hecke function has no pole.
On every fixed closed vertical strip, \(\mathcal MF\)
and its derivatives decay faster than every power of
\(|\operatorname{Im}z|\), whereas the other factors and
their derivatives grow at most polynomially. Thus the
same rapid bounds hold for \(H_F\) and all its derivatives.
These transform facts follow directly by integration by
parts for \(F\in\mathcal E\), and by the fixed-strip
Hecke bounds just stated.

Inverse Mellin integration along any real line defines
the same smooth function: shift a finite-height rectangle
between two lines and let its height tend to infinity.
The horizontal integrals tend to zero by these uniform
rapid bounds, and no residue is crossed. On each chosen
line its exponentially weighted inverse is Schwartz,
by Fourier inversion. Since the line is arbitrary, that
inverse belongs to \(\mathcal E\). On a line \(b>1/2\)
it is the literal sum just proved, so the two are equal
everywhere by smoothness. Hence
\[
 \mathcal M(Z_{\eta,S}^{\rm geom}F)(z)=H_F(z)
                         \quad\hbox{for every }z\in\mathbb C.
 \tag{MA-MO.110}
\]
This proves actual equality on the narrow line, not merely
a definition of its output by analytic continuation.

At every completed nontrivial zero \(\rho\), the finite
\(L_\eta\) has at least the completed multiplicity: the
gamma factors have no zeros. Thus (MA-MO.110) vanishes to
full required order at every \(\rho-1/2\). The exact
kernel equality (M18.4), with the explicit center
\(T_{1/2}F=e^{x/2}F\), gives
\(Z_{\eta,S}^{\rm geom}F\in T_{1/2}\mathcal I_\eta\).
The actual theta-image equality (M24.9), together with
(M24.3), identifies this with the original smooth
compact-character theta image. This argument uses the
proved image equality and does not invent an adelic
radial representative for each ideal class.

The original wide inverse of (MA-CC.1) returns \(F\),
since its absolutely convergent ideal M\"obius series is
inverse to \(Z\) on \(\mathcal H_b\), \(b>1/2\).
On the narrow line, (MA-MO.110) gives
\(V_ay=A V_aF\). Since \(F\in\mathcal H_a\) and
\(jA=1\) almost everywhere, \(y\) is in the maximal
inverse domain and its inverse is again \(F\).
The same argument applies to \(\phi F\in\mathcal E\).
It proves every equality and domain assertion in
(MA-MO.104), with no moment conditions.

\emph{Graph approximation by the actual smooth ideal.}
Fix \(v\in\mathscr C_{\eta,\sigma}\), put \(q_0=jv\),
and choose \(h_0=\mathcal F_+^{-1}q_0\in\mathscr S\).
For \(\epsilon>0\) set
\[
 F_\epsilon(x)=e^{-\epsilon x^2-ax}h_0(x),\qquad
 q_\epsilon=V_aF_\epsilon
             =\mathcal F_+(e^{-\epsilon x^2}h_0).
 \tag{MA-MO.111}
\]
Every \(F_\epsilon\) belongs to \(\mathcal E\), and
\(q_\epsilon\to q_0\) in Schwartz space. The latter
follows by the product rule: Gaussian derivatives are
uniformly polynomially bounded for \(0<\epsilon\le1\),
their derivatives converge on compact sets, and the
Schwartz tails of \(h_0\) control the complementary
region in every seminorm. Fourier transformation
preserves that convergence.

Let \(y_\epsilon=Z_{\eta,S}^{\rm geom}F_\epsilon\).
The equalities already proved and polynomial bounds
on \(A\) give
\[
 \begin{aligned}
 V_ay_\epsilon&=Aq_\epsilon\longrightarrow v,\\
 V_aQ_{\eta,\phi,S}y_\epsilon
       &=AT_\phi q_\epsilon
                    \longrightarrow AT_\phi q_0
                         \qquad\text{in }L^2.
 \end{aligned}
 \tag{MA-MO.112}
\]
Continuity of \(T_\phi\) on Schwartz space justifies
the second limit. Every vector and every cutoff output
in this display is the literal smooth ideal sum of
(MA-MO.104). Thus the graph closure of this restriction
contains the entire compact-frequency core graph.

For an explicit unbounded sequence, use the normalized
core vectors from (MA-MO.108), and for each \(M\) choose
\(\epsilon_M\) so that the input and output errors in
(MA-MO.112) are at most one. Then
\(\|y_{\epsilon_M}\|_{\mathcal H_a}\le2\), whereas
\(\|Q_{\eta,\phi,S}y_{\epsilon_M}\|_{\mathcal H_a}
\ge c_S e^{Md}/M-1\to\infty\).
The bounded multiplier \(M_\phi\) gives the same
divergence for the difference. This proves the claimed
sequence and the literal geometric unboundedness.

All preceding smooth-domain steps use only that the
fixed nontrivial Hecke function is entire with the
standard strip bounds. For a primitive nontrivial
Dirichlet function these hypotheses hold as well,
and MA-MO-002 supplies the first norm obstruction.
For the conjugate imaginary-quadratic character,
Mishou's theorem and each argument apply unchanged.
The reflected source component uses that conjugate
character exactly as in (MA-CC.9) and (MA-RS.44).
\end{proof}

\paragraph{Original source scope and zero-location bookkeeping.}
For the source cutoff and a structural place set \(S\) as in MA-CC-001,
(MA-CC.8) identifies the literal ideal operator with
the actual compact-character restriction of Meyer's
Euler summation. Hence the source first diagonal
\(M_\phi-Z_{\eta,S}^{\rm geom}M_\phi
J_{\eta,S}^{\rm geom}\) has the stated obstruction
on genuine smooth ideal inputs. This refutes a bounded
extension of that diagonal, and of a full product
matrix which agrees on its coordinate core.
It does not rule out cancellations on a separately
restricted common-sum domain or a smoothed operator.
It does not change the discarded-divisor gate in
MA-FT-002 or prove Hecke GRH for the cubic example.

If a zero is assumed to lie on the chosen line,
the analytic multiplier \(A\) satisfies every
hypothesis of MA-MO-006: it is analytic near the
line zero, its finite factors there are nonzero,
and it and its derivatives have polynomial growth.
That already proved theorem therefore also gives
conditional nonclosability on the compact-frequency
core, even after any nonzero compact left smoothing.
No existence of such a zero is asserted, and this
core conclusion is not transferred without proof
to the distinct geometric restriction above.

In Grothendieck's inherited trace setting, the
actual topological operator remains the required
input before a nuclear trace or Fredholm determinant
can be used: cote~1, scans~079--080 and~111--112,
through ALPOGE\_BRIDGES BR-TR-01/02. The character
is the rank-one finite-order local system already
identified in MA-CH/MA-CQ; prime-ideal coefficients
replace rational primes in its Euler realization.
No new weight, positivity, or motivic assertion is
being substituted for the analytic domain check.

Mishou, H. (2003). The universality theorem for Hecke
\(L\)-functions. \emph{Acta Arithmetica, 110}(1), 45--71.
\url{https://doi.org/10.4064/aa110-1-4}.
Exact locators: main unnumbered Theorem, p.~46;
Gr\"ossencharacter definition, \S2, pp.~47--49.
The primary statement and its applicability were freshly
read. Its complete classical proof is carried as
published background, not recorded as a new full proof
audit. The new operator deduction above receives its
own complete reading events.
